\chapter{Model Correction}
\label{chap:ModelCorrection}

\section{Order-Dependent Characteristics}
\label{sec:order_dependent_characteristics}

An intuitive approach to solve the order-dependence of the microcluster density is to redefine it as an aggregation of the Typicality and Eccentricity Data Analytics variables across all aggregated samples. We show, however, that such solution is analytically flawed.

If

\begin{equation}
    V_k = \sum_{j=1}^{S_k} \lVert \mu_k - x_j \rVert^2
\end{equation}

is the sum of squared Euclidean deviations for the $S_k$ samples of a microcluster $i$, for the standard spatial sample variance $\sigma^2_k = \frac{1}{S_k} V_k$,

\begin{equation}
    V_k = S_k \sigma^2_k
\end{equation}

Then, $\bar{\xi}_k$, the mean eccentricity of a microcluster $k$, is defined as \cref{eq:mean_eccentricity}.

\begin{align}
    \bar{\xi}_k &= \frac{1}{S_k} \sum_{j=1}^{S_k} \left( \frac{1}{S_k} + \frac{\lVert \mu_k - x_j \rVert^2}{S_k \sigma^2_k} \right) \nonumber \\
                &= \frac{1}{S_k} \left( \sum_{j=1}^{S_k} \frac{1}{S_k} + \frac{1}{S_k \sigma^2_k} \sum_{j=1}^{S_k} \lVert \mu_k - x_j \rVert^2 \right) \nonumber \\
                &= \frac{1}{S_k} \left( 1 + \frac{V_k}{S_k \sigma^2_k} \right) \nonumber \\
                &= \frac{1}{S_k} \left( 1 + \frac{S_k \sigma^2_k}{S_k \sigma^2_k} \right) = \frac{2}{S_k}
    \label{eq:mean_eccentricity}
\end{align}

If the exact recursive formulas from the original \MicroTEDAClus formulation are injected instead \cite{maia2020evolvclus}, the cancellation persists. The original formulation incorporates an anomalous penalty by the dimensionality $d$ squared in its variance computation, such that

\begin{equation}
    \sigma^2_k = \frac{1}{S_k} \sum_{j=1}^{S_k} \frac{4\lVert \mu_k - x_j \rVert^2}{d^2}
\end{equation}

Which corresponts to $V_k = \frac{d^2 S_k \sigma^2_k}{4}$. The original eccentricity formula is \cref{eq:original_eccentricity}.

\begin{equation}
    \xi(x_j) = \frac{1}{S_k} + \frac{2\lVert \mu_k - x_j \rVert^2}{S_k \sigma^2_k d}
    \label{eq:original_eccentricity}
\end{equation}

Thus,

\begin{align}
    \bar{\xi}_k &= \frac{1}{S_k} \sum_{j=1}^{S_k} \left( \frac{1}{S_k} + \frac{2\lVert \mu_k - x_j \rVert^2}{S_k \sigma^2_k d} \right) \nonumber \\
                &= \frac{1}{S_k} \left( 1 + \frac{2}{S_k \sigma^2_k d} \sum_{j=1}^{S_k} \lVert \mu_k - x_j \rVert^2 \right) \nonumber \\
                &= \frac{1}{S_k} \left( 1 + \frac{2 V_k}{S_k \sigma^2_k d} \right) \nonumber \\
                &= \frac{1}{S_k} \left( 1 + \frac{2}{S_k \sigma^2_k d} \left( \frac{d^2 S_k \sigma^2_k}{4} \right) \right) \nonumber \\
                &= \frac{1}{S_k} \left( 1 + \frac{d}{2} \right)
\end{align}

In both cases, the spatial dispersion term $\sigma^2_k$ cancels out entirely. Because the geometric volume is analytically eliminated, the aggregated TEDA variables reduce strictly to a function of cardinality ($S_k$) and dimensionality ($d$). Consequently, any density metric derived from them collapses into a pure mass count, rendering the algorithm mathematically blind to actual spatial dispersion.

To formulate a geometrically valid and order-invariant density, the microcluster characterization is partitioned into intrinsic state variables and transient relational metrics. The TEDA variables $\xi(x)$, $\zeta(x)$, $\tau(x)$, and $t(x)$ are strictly preserved as point-to-set transient functions evaluated for an incoming sample $x$ against the distribution \cite{angelov2014outside}. They are explicitly removed from the microcluster's intrinsic state variables.

The microcluster state is thus restricted to its recursively updated sufficient statistics: $S_k$, $\mu_k$, and $\sigma^2_k$. Because the model characterizes the spatial dispersion of a $d$-dimensional microcluster using a single scalar variance $\sigma^2_k$ rather than a full covariance matrix, it mathematically enforces an isotropic geometric assumption. Under this assumption, the true spatial volume scales with $(\sigma^2_k)^{d/2}$. To avoid the numerical instability of exponential dimensionality scaling in continuous stream clustering, a linear proxy of variance is utilized as a standard algebraic relaxation \cite{cao2006density}. The microcluster density $D_k$ is redefined directly as the ratio of accumulated mass to this volumetric proxy in \cref{eq:corrected_density}.

\begin{equation}
    D_k = \frac{S_k}{\sigma^2_k + \epsilon}
    \label{eq:corrected_density}
\end{equation}

where $\epsilon > 0$ is a stabilization constant to prevent undefined behavior in single-sample microclusters.

\Cref{tab:orderdep_ari_overall,tab:orderdep_ari_fifth,tab:orderdep_ari_steadystate,tab:orderdep_ari_diff} present the clustering quality impact of the correction. Although it harmed the overall and initial ARI for four of the six datasets, it did not significantly affect the steady-state ARI measures, even improving them for the Chameleon sets. Our hypothesis is that the original order-dependent density formulation better match the remaining flaws in the model, thus compensating for them. This is investigated through the following sections, in which the model is further corrected.

\begin{table}[htbp]
\centering
\caption{Statistical analysis and (Q1, Q3) for the Overall Adjusted Rand Index.}
\resizebox{\textwidth}{!}{%
\begin{tabular}{l c c }
\toprule
Dataset & \modelname{MicroTEDAClus} & \modelname{MicroTEDAClusCorrected} \\
\midrule
\datasetname{Stream Benchmark 1} & \textbf{0.019 (0.000, 0.453)} & \textcolor{red}{0.000 (0.000, 0.021)} \\
\datasetname{Stream Benchmark 2} & \textbf{0.577 (0.517, 0.626)} & \textcolor{red}{0.503 (0.223, 0.547)} \\
\datasetname{Random RBF Generator with Events} & \textbf{0.459 (0.140, 0.588)} & \textcolor{red}{0.000 (0.000, 0.343)} \\
\datasetname{Cassini} & \textbf{0.456 (0.066, 0.547)} & \textcolor{red}{0.220 (0.000, 0.352)} \\
\datasetname{Chameleon DS3} & 0.208 (0.150, 0.240) & 0.195 (0.149, 0.242) \\
\datasetname{Chameleon DS4} & \textcolor{red}{0.160 (0.125, 0.223)} & \textbf{0.207 (0.161, 0.235)} \\
\bottomrule
\end{tabular}
}
\label{tab:orderdep_ari_overall}
\end{table}

\begin{table}[htbp]
\centering
\caption{Statistical analysis for the Adjusted Rand Index at the Fifth Batch.}
\resizebox{\textwidth}{!}{%
\begin{tabular}{l c c }
\toprule
Dataset & \modelname{baseline} & \modelname{order-independent} \\
\midrule
\datasetname{Stream Benchmark 1} & \textbf{0.852 (0.802, 0.883)} & \textcolor{red}{0.782 (0.000, 0.851)} \\
\datasetname{Stream Benchmark 2} & \textbf{0.591 (0.552, 0.616)} & \textcolor{red}{0.542 (0.217, 0.598)} \\
\datasetname{Random RBF Generator with Events} & \textbf{0.560 (0.403, 0.802)} & \textcolor{red}{0.213 (0.055, 0.397)} \\
\datasetname{Cassini} & \textbf{0.502 (0.360, 0.561)} & \textcolor{red}{0.207 (0.000, 0.420)} \\
\datasetname{Chameleon DS3} & 0.253 (0.144, 0.255) & 0.218 (0.137, 0.279) \\
\datasetname{Chameleon DS4} & 0.102 (0.000, 0.188) & 0.102 (0.000, 0.174) \\
\bottomrule
\end{tabular}
}
\label{tab:orderdep_ari_fifth}
\end{table}

\begin{table}[htbp]
\centering
\caption{Statistical analysis for the Steady-State Adjusted Rand Index.}
\resizebox{\textwidth}{!}{%
\begin{tabular}{l c c }
\toprule
Dataset & \modelname{baseline} & \modelname{order-independent} \\
\midrule
\datasetname{Stream Benchmark 1} & 0.000 (0.000, 0.025) & 0.007 (0.000, 0.279) \\
\datasetname{Stream Benchmark 2} & 0.538 (0.076, 0.645) & 0.569 (0.505, 0.650) \\
\datasetname{Random RBF Generator with Events} & \textcolor{red}{0.054 (0.000, 0.453)} & \textbf{0.403 (0.187, 0.554)} \\
\datasetname{Cassini} & \textbf{0.458 (0.174, 0.557)} & \textcolor{red}{0.288 (0.186, 0.356)} \\
\datasetname{Chameleon DS3} & 0.187 (0.162, 0.194) & 0.188 (0.171, 0.200) \\
\datasetname{Chameleon DS4} & 0.217 (0.178, 0.247) & 0.237 (0.169, 0.263) \\
\datasetname{SCL-32} & \textbf{0.843 (0.000, 0.955)} & \textcolor{red}{0.029 (0.000, 0.873)} \\
\datasetname{SCL-64} & 0.312 (0.000, 0.864) & 0.000 (0.000, 0.845) \\
\bottomrule
\end{tabular}
}
\label{tab:orderdep_ari_steadystate}
\end{table}

\begin{table}[htbp]
\centering
\caption{Statistical analysis for the Maximum ARI Drop.}
\resizebox{\textwidth}{!}{%
\begin{tabular}{l c c }
\toprule
Dataset & \modelname{baseline} & \modelname{order-independent} \\
\midrule
\datasetname{Stream Benchmark 1} & 0.921 (0.902, 0.959) & 0.907 (0.869, 0.963) \\
\datasetname{Stream Benchmark 2} & 0.895 (0.665, 1.000) & 0.929 (0.759, 0.971) \\
\datasetname{Random RBF Generator with Events} & 0.900 (0.698, 0.961) & 0.901 (0.806, 0.975) \\
\datasetname{Cassini} & \textcolor{red}{0.743 (0.652, 0.791)} & \textbf{0.598 (0.465, 0.713)} \\
\datasetname{Chameleon DS3} & 0.426 (0.426, 0.426) & 0.426 (0.402, 0.427) \\
\datasetname{Chameleon DS4} & 0.442 (0.380, 0.484) & 0.400 (0.380, 0.455) \\
\datasetname{SCL-32} & 1.000 (0.947, 1.000) & 1.000 (1.000, 1.002) \\
\datasetname{SCL-64} & 1.000 (0.951, 1.002) & 1.000 (1.000, 1.003) \\
\bottomrule
\end{tabular}
}
\label{tab:orderdep_ari_diff}
\end{table}


Regarding time efficiency and resource consumption, \cref{tab:orderdep_time_fit_cum_ms,tab:orderdep_time_predict_cum_ms,tab:orderdep_memory_max_bytes,tab:orderdep_memory_steadystate_bytes,tab:orderdep_mc_max,tab:orderdep_mc_steadystate} show that the correction provided an overall improvement, specially for the datasets that feature concept drift. Under all those metrics, the corrected model is either better or equivalent to the original formulation.

\begin{table}[htbp]
\centering
\caption{Statistical analysis for the Cumulative Fitting Time [ms].}
\resizebox{\textwidth}{!}{%
\begin{tabular}{l c c }
\toprule
Dataset & \modelname{baseline} & \modelname{order-independent} \\
\midrule
\datasetname{Stream Benchmark 1} & 440.0 (200.0, 1675.0) & 1150.0 (557.5, 1560.0) \\
\datasetname{Stream Benchmark 2} & \textcolor{red}{555.0 (97.5, 927.5)} & \textbf{280.0 (175.0, 395.0)} \\
\datasetname{Random RBF Generator with Events} & \textcolor{red}{6940.0 (2052.5, 18825.0)} & \textbf{1760.0 (932.5, 3010.0)} \\
\datasetname{Cassini} & \textbf{670.0 (600.0, 865.0)} & \textcolor{red}{1380.0 (390.0, 1955.0)} \\
\datasetname{Chameleon DS3} & 5350.0 (5250.0, 12250.0) & 5605.0 (4092.5, 33760.0) \\
\datasetname{Chameleon DS4} & 5315.0 (4530.0, 6785.0) & 5005.0 (4235.0, 8592.5) \\
\bottomrule
\end{tabular}
}
\label{tab:orderdep_time_fit_cum_ms}
\end{table}

\begin{table}[htbp]
\centering
\caption{Statistical analysis for the Cumulative Prediction Time [ms].}
\resizebox{\textwidth}{!}{%
\begin{tabular}{l c c }
\toprule
Dataset & \modelname{baseline} & \modelname{order-independent} \\
\midrule
\datasetname{Stream Benchmark 1} & 50.0 (20.0, 145.0) & 75.0 (47.5, 110.0) \\
\datasetname{Stream Benchmark 2} & \textcolor{red}{60.0 (20.0, 80.0)} & \textbf{30.0 (20.0, 42.5)} \\
\datasetname{Random RBF Generator with Events} & \textcolor{red}{845.0 (197.5, 3252.5)} & \textbf{135.0 (85.0, 240.0)} \\
\datasetname{Cassini} & \textcolor{red}{70.0 (60.0, 82.5)} & \textbf{40.0 (30.0, 50.0)} \\
\datasetname{Chameleon DS3} & \textcolor{red}{490.0 (450.0, 795.0)} & \textbf{170.0 (160.0, 585.0)} \\
\datasetname{Chameleon DS4} & \textcolor{red}{280.0 (197.5, 430.0)} & \textbf{115.0 (90.0, 230.0)} \\
\bottomrule
\end{tabular}
}
\label{tab:orderdep_time_predict_cum_ms}
\end{table}

\begin{table}[htbp]
\centering
\caption{Statistical analysis for the Maximum Memory Footprint [B].}
\resizebox{\textwidth}{!}{%
\begin{tabular}{l c c }
\toprule
Dataset & \modelname{baseline} & \modelname{order-independent} \\
\midrule
\datasetname{Stream Benchmark 1} & \textbf{151040 (88576, 329984)} & \textcolor{red}{310784 (162816, 489472)} \\
\datasetname{Stream Benchmark 2} & \textcolor{red}{146944 (78592, 159488)} & \textbf{80384 (43264, 90112)} \\
\datasetname{Random RBF Generator with Events} & \textcolor{red}{834560 (268800, 1473280)} & \textbf{299520 (214784, 425984)} \\
\datasetname{Cassini} & 144896 (140288, 152320) & 158720 (129536, 183040) \\
\datasetname{Chameleon DS3} & 325120 (318976, 470784) & 282624 (274432, 873728) \\
\datasetname{Chameleon DS4} & \textcolor{red}{279552 (270848, 334848)} & \textbf{252928 (245760, 367616)} \\
\bottomrule
\end{tabular}
}
\label{tab:orderdep_memory_max_bytes}
\end{table}

\begin{table}[htbp]
\centering
\caption{Statistical analysis for the Steady-State Memory Footprint [B].}
\resizebox{\textwidth}{!}{%
\begin{tabular}{l c c }
\toprule
Dataset & \modelname{baseline} & \modelname{order-independent} \\
\midrule
\datasetname{Stream Benchmark 1} & 67135 (17499, 322840) & 188480 (115516, 515336) \\
\datasetname{Stream Benchmark 2} & 90504 (11791, 136409) & 120805 (35672, 210991) \\
\datasetname{Random RBF Generator with Events} & 307390 (188323, 1218522) & 557951 (242417, 1366732) \\
\datasetname{Cassini} & 123210 (113289, 131625) & 102881 (49032, 118929) \\
\datasetname{Chameleon DS3} & 317525 (314250, 489585) & 307577 (305326, 431567) \\
\datasetname{Chameleon DS4} & 280885 (266411, 400677) & 264512 (256289, 395176) \\
\datasetname{SCL-32} & 19476 (14216, 704242) & 19596 (14637, 107708) \\
\datasetname{SCL-64} & 19579 (15685, 2292394) & 18014 (15832, 100725) \\
\bottomrule
\end{tabular}
}
\label{tab:orderdep_memory_steadystate_bytes}
\end{table}

\begin{table}[htbp]
\centering
\caption{Statistical analysis for the Maximum Microcluster Count.}
\resizebox{\textwidth}{!}{%
\begin{tabular}{l c c }
\toprule
Dataset & \modelname{baseline} & \modelname{order-independent} \\
\midrule
\datasetname{Stream Benchmark 1} & 36 (15, 130) & 60 (41, 84) \\
\datasetname{Stream Benchmark 2} & 44 (14, 62) & 42 (28, 82) \\
\datasetname{Random RBF Generator with Events} & 228 (57, 432) & 131 (60, 396) \\
\datasetname{Cassini} & 46 (40, 56) & 40 (27, 46) \\
\datasetname{Chameleon DS3} & 128 (124, 284) & 124 (119, 226) \\
\datasetname{Chameleon DS4} & 124 (116, 224) & 113 (106, 221) \\
\datasetname{SCL-32} & 35 (17, 79) & 28 (20, 60) \\
\datasetname{SCL-64} & 32 (26, 89) & 26 (25, 48) \\
\bottomrule
\end{tabular}
}
\label{tab:orderdep_mc_max}
\end{table}

\begin{table}[htbp]
\centering
\caption{Statistical analysis for the Steady-State Microcluster Count.}
\resizebox{\textwidth}{!}{%
\begin{tabular}{l c c }
\toprule
Dataset & \modelname{baseline} & \modelname{order-independent} \\
\midrule
\datasetname{Stream Benchmark 1} & 13 (4, 106) & 47 (30, 78) \\
\datasetname{Stream Benchmark 2} & 24 (6, 46) & 35 (16, 67) \\
\datasetname{Random RBF Generator with Events} & 71 (15, 342) & 115 (42, 296) \\
\datasetname{Cassini} & 39 (32, 45) & 34 (18, 42) \\
\datasetname{Chameleon DS3} & 120 (118, 271) & 115 (113, 217) \\
\datasetname{Chameleon DS4} & 119 (110, 218) & 108 (101, 214) \\
\datasetname{SCL-32} & 7 (4, 35) & 8 (5, 23) \\
\datasetname{SCL-64} & 6 (5, 22) & 5 (4, 25) \\
\bottomrule
\end{tabular}
}
\label{tab:orderdep_mc_steadystate}
\end{table}


\section{Dimensionality Penalty}
\label{sec:dimensionality_penalty}

In the original formulation \cite{maia2020evolvclus}, the spatial dispersion metrics explicitly penalize the Euclidean norm by the data dimensionality $d$. Specifically, the recursive variance $\left(\sigma^i_k\right)^2$ and the point-to-set eccentricity $\xi^i(x_k)$ arbitrarily scale the squared deviations by $4/d^2$ and $2/d$, respectively. This scaling factor is analytically flawed. It distorts the true geometric boundaries of the microclusters and deviates from the standard TEDA derivations for the Euclidean space \cite{angelov2014outside}, where the spatial variance strictly reflects the unscaled trace of the covariance matrix.

To restore the analytical integrity of the spatial boundaries, the dimensionality penalty is removed. The recursive variance update for a microcluster $i$ at moment $k$ is corrected to the standard formulation:

$$ \left(\sigma^i_k\right)^2 = \frac{S^i_k - 1}{S^i_k} \left(\sigma^i_{k-1}\right)^2 + \frac{1}{S^i_k - 1} \lVert x_k - \mu^i_k \rVert^2 $$

Correspondingly, the transient point-to-set eccentricity evaluation is restored to its exact unscaled geometric form:

$$ \xi^i(x_k) = \frac{1}{S^i_k} + \frac{\lVert \mu^i_k - x_k \rVert^2}{S^i_k \left(\sigma^i_k\right)^2} $$

After this correction, the original model is shown by \cref{tab:dimpen_ari_overall,tab:dimpen_ari_fifth} to continue the best implementation under the overall and initial ARI measures. Nevertheless, \cref{tab:dimpen_ari_steadystate,tab:dimpen_ari_diff} show a pronounced improvement in the steady-state and ARI drop, mostly for the higher dimensional datasets.

\begin{table}[htbp]
\centering
\caption{Statistical analysis for the Overall Adjusted Rand Index.}
\resizebox{\textwidth}{!}{%
\begin{tabular}{l c c c }
\toprule
Dataset & \modelname{baseline} & \modelname{order-independent} & \modelname{no-dim-penalty} \\
\midrule
\datasetname{Stream Benchmark 1} & 0.368 (0.283, 0.416) & 0.286 (0.155, 0.386) & 0.312 (0.155, 0.381) \\
\datasetname{Stream Benchmark 2} & 0.584 (0.364, 0.617) & 0.578 (0.507, 0.610) & 0.567 (0.490, 0.613) \\
\datasetname{Random RBF Generator with Events} & 0.356 (0.142, 0.518) & 0.418 (0.349, 0.485) & 0.412 (0.331, 0.485) \\
\datasetname{Cassini} & \textbf{0.476 (0.329, 0.509)} & \textcolor{red}{0.312 (0.183, 0.349)} & \textcolor{red}{0.289 (0.111, 0.349)} \\
\datasetname{Chameleon DS3} & 0.228 (0.200, 0.238) & 0.232 (0.206, 0.238) & 0.233 (0.207, 0.241) \\
\datasetname{Chameleon DS4} & 0.242 (0.219, 0.263) & 0.259 (0.219, 0.273) & 0.262 (0.219, 0.273) \\
\datasetname{SCL-32} & \textcolor{red}{0.828 (0.398, 0.884)} & \textcolor{red}{0.710 (0.415, 0.861)} & \textbf{0.836 (0.833, 0.837)} \\
\datasetname{SCL-64} & \textcolor{red}{0.775 (0.561, 0.831)} & \textcolor{red}{0.661 (0.505, 0.829)} & \textbf{0.836 (0.833, 0.838)} \\
\bottomrule
\end{tabular}
}
\label{tab:dimpen_ari_overall}
\end{table}

\begin{table}[htbp]
\centering
\caption{Statistical analysis for the Adjusted Rand Index at the Fifth Batch.}
\resizebox{\textwidth}{!}{%
\begin{tabular}{l c c c }
\toprule
Dataset & \modelname{baseline} & \modelname{order-independent} & \modelname{no-dim-penalty} \\
\midrule
\datasetname{Stream Benchmark 1} & \textbf{0.852 (0.802, 0.883)} & \textcolor{red}{0.782 (0.000, 0.851)} & \textcolor{red}{0.734 (0.000, 0.841)} \\
\datasetname{Stream Benchmark 2} & 0.591 (0.552, 0.616) & 0.542 (0.217, 0.598) & 0.564 (0.372, 0.605) \\
\datasetname{Random RBF Generator with Events} & \textbf{0.560 (0.403, 0.802)} & \textcolor{red}{0.213 (0.055, 0.397)} & \textcolor{red}{0.154 (0.041, 0.348)} \\
\datasetname{Cassini} & \textbf{0.502 (0.360, 0.561)} & \textcolor{red}{0.207 (0.000, 0.420)} & \textcolor{red}{0.045 (0.000, 0.402)} \\
\datasetname{Chameleon DS3} & 0.253 (0.144, 0.255) & 0.218 (0.137, 0.279) & 0.218 (0.147, 0.279) \\
\datasetname{Chameleon DS4} & 0.102 (0.000, 0.188) & 0.102 (0.000, 0.174) & 0.102 (0.000, 0.174) \\
\datasetname{SCL-32} & 0.903 (0.841, 0.997) & - & 0.879 (0.856, 0.905) \\
\datasetname{SCL-64} & \textbf{0.903 (0.859, 0.972)} & - & \textcolor{red}{0.867 (0.828, 0.893)} \\
\bottomrule
\end{tabular}
}
\label{tab:dimpen_ari_fifth}
\end{table}

\begin{table}[htbp]
\centering
\caption{Statistical analysis and (Q1, Q3) for the Adjusted Rand Index at the Penultimate Batch.}
\resizebox{\textwidth}{!}{%
\begin{tabular}{l c c c }
\toprule
Dataset & \modelname{baseline} & \modelname{order-independent} & \modelname{no-dim-penalty} \\
\midrule
\datasetname{Stream Benchmark 1} & 0.000 (0.000, 0.037) & 0.000 (0.000, 0.035) & 0.000 (0.000, 0.055) \\
\datasetname{Stream Benchmark 2} & 0.491 (0.000, 0.768) & 0.428 (0.020, 0.535) & 0.372 (0.000, 0.500) \\
\datasetname{Random RBF Generator with Events} & 0.076 (0.000, 0.451) & 0.003 (0.000, 0.415) & 0.007 (0.000, 0.348) \\
\datasetname{Cassini} & 0.396 (0.029, 0.581) & 0.306 (0.066, 0.437) & 0.317 (0.002, 0.410) \\
\datasetname{Chameleon DS3} & \textcolor{red}{0.129 (0.057, 0.219)} & 0.157 (0.069, 0.229) & \textbf{0.207 (0.166, 0.241)} \\
\datasetname{Chameleon DS4} & \textcolor{red}{0.098 (0.068, 0.153)} & 0.146 (0.114, 0.190) & \textbf{0.266 (0.231, 0.281)} \\
\datasetname{SCL-32} & \textcolor{red}{0.662 (0.000, 0.924)} & - & \textbf{0.875 (0.845, 0.905)} \\
\datasetname{SCL-64} & \textcolor{red}{0.000 (0.000, 0.843)} & - & \textbf{0.863 (0.828, 0.897)} \\
\bottomrule
\end{tabular}
}
\label{tab:dimpen_ari_steadystate}
\end{table}

\begin{table}[htbp]
\centering
\caption{Statistical analysis for the Maximum ARI Drop.}
\resizebox{\textwidth}{!}{%
\begin{tabular}{l c c c }
\toprule
Dataset & \modelname{baseline} & \modelname{order-independent} & \modelname{no-dim-penalty} \\
\midrule
\datasetname{Stream Benchmark 1} & 0.921 (0.902, 0.959) & 0.907 (0.869, 0.963) & 0.920 (0.881, 0.953) \\
\datasetname{Stream Benchmark 2} & 0.895 (0.665, 1.000) & 0.929 (0.759, 0.971) & 0.916 (0.693, 0.970) \\
\datasetname{Random RBF Generator with Events} & 0.900 (0.698, 0.961) & 0.901 (0.806, 0.975) & 0.901 (0.806, 0.974) \\
\datasetname{Cassini} & \textcolor{red}{0.743 (0.652, 0.791)} & \textbf{0.598 (0.465, 0.713)} & \textbf{0.592 (0.462, 0.722)} \\
\datasetname{Chameleon DS3} & 0.426 (0.426, 0.426) & 0.426 (0.402, 0.427) & 0.426 (0.401, 0.438) \\
\datasetname{Chameleon DS4} & 0.442 (0.380, 0.484) & 0.400 (0.380, 0.455) & 0.400 (0.380, 0.455) \\
\datasetname{SCL-32} & 1.000 (0.947, 1.000) & \textcolor{red}{1.000 (1.000, 1.002)} & \textbf{0.940 (0.935, 0.953)} \\
\datasetname{SCL-64} & \textcolor{red}{1.000 (0.951, 1.002)} & \textcolor{red}{1.000 (1.000, 1.003)} & \textbf{0.942 (0.937, 0.948)} \\
\bottomrule
\end{tabular}
}
\label{tab:dimpen_ari_diff}
\end{table}


\Cref{tab:dimpen_time_fit_cum_ms,tab:dimpen_time_predict_cum_ms,tab:dimpen_memory_max_bytes,tab:dimpen_memory_steadystate_bytes,tab:dimpen_mc_max,tab:dimpen_mc_steadystate} show that the improvement achieved by the previous correction is preserved for the first set of datasets. For the newly included \datasetname{SCL-32} and \datasetname{SCL-64}, the scores are consistently worst. This is expected, as thoroughly discussed in \cref{sec:reassessment_results}, since the dimensionality penalty colapses the spatial structure into fewer, lower quality microclusters, which are more efficient to maintain and evaluate, but less representative of the underlying distribution. The corrected model, which preserves the true geometric boundaries, is more computationally expensive, but it is also more accurate for those datasets.

\begin{table}[htbp]
\centering
\caption{Statistical analysis for the Cumulative Fitting Time [ms].}
\resizebox{\textwidth}{!}{%
\begin{tabular}{l c c c }
\toprule
Dataset & \modelname{baseline} & \modelname{order-independent} & \modelname{no-dim-penalty} \\
\midrule
\datasetname{Stream Benchmark 1} & 440.0 (200.0, 1675.0) & 1150.0 (557.5, 1560.0) & 975.0 (530.0, 1422.5) \\
\datasetname{Stream Benchmark 2} & 555.0 (97.5, 927.5) & 280.0 (175.0, 395.0) & 240.0 (167.5, 397.5) \\
\datasetname{Random RBF Generator with Events} & 6940.0 (2052.5, 18825.0) & 1760.0 (932.5, 3010.0) & 1430.0 (932.5, 3215.0) \\
\datasetname{Cassini} & 670.0 (600.0, 865.0) & 1380.0 (390.0, 1955.0) & 1390.0 (397.5, 2202.5) \\
\datasetname{Chameleon DS3} & 5350.0 (5250.0, 12250.0) & 5605.0 (4092.5, 33760.0) & 5580.0 (4115.0, 33670.0) \\
\datasetname{Chameleon DS4} & 5315.0 (4530.0, 6785.0) & 5005.0 (4235.0, 8592.5) & 4995.0 (4292.5, 8507.5) \\
\datasetname{SCL-32} & 515.0 (232.5, 1440.0) & - & 805.0 (457.5, 1680.0) \\
\datasetname{SCL-64} & \textbf{335.0 (260.0, 967.5)} & - & \textcolor{red}{1230.0 (495.0, 3207.5)} \\
\bottomrule
\end{tabular}
}
\label{tab:dimpen_time_fit_cum_ms}
\end{table}

\begin{table}[htbp]
\centering
\caption{Statistical analysis for the Cumulative Prediction Time [ms].}
\resizebox{\textwidth}{!}{%
\begin{tabular}{l c c c }
\toprule
Dataset & \modelname{baseline} & \modelname{order-independent} & \modelname{no-dim-penalty} \\
\midrule
\datasetname{Stream Benchmark 1} & 50.0 (20.0, 145.0) & 75.0 (47.5, 110.0) & 70.0 (40.0, 85.0) \\
\datasetname{Stream Benchmark 2} & 60.0 (20.0, 80.0) & 30.0 (20.0, 42.5) & 30.0 (20.0, 40.0) \\
\datasetname{Random RBF Generator with Events} & \textcolor{red}{845.0 (197.5, 3252.5)} & \textbf{135.0 (85.0, 240.0)} & \textbf{130.0 (70.0, 272.5)} \\
\datasetname{Cassini} & \textcolor{red}{70.0 (60.0, 82.5)} & \textbf{40.0 (30.0, 50.0)} & \textbf{40.0 (30.0, 40.0)} \\
\datasetname{Chameleon DS3} & \textcolor{red}{490.0 (450.0, 795.0)} & \textbf{170.0 (160.0, 585.0)} & \textbf{165.0 (150.0, 572.5)} \\
\datasetname{Chameleon DS4} & \textcolor{red}{280.0 (197.5, 430.0)} & \textbf{115.0 (90.0, 230.0)} & \textbf{110.0 (87.5, 235.0)} \\
\datasetname{SCL-32} & \textcolor{red}{65.0 (47.5, 117.5)} & - & \textbf{40.0 (30.0, 55.0)} \\
\datasetname{SCL-64} & \textcolor{red}{65.0 (50.0, 90.0)} & - & \textbf{40.0 (30.0, 70.0)} \\
\bottomrule
\end{tabular}
}
\label{tab:dimpen_time_predict_cum_ms}
\end{table}

\begin{table}[htbp]
\centering
\caption{Statistical analysis and (Q1, Q3) for the Maximum Memory Footprint [B].}
\resizebox{\textwidth}{!}{%
\begin{tabular}{l c c c }
\toprule
Dataset & \modelname{baseline} & \modelname{order-independent} & \modelname{no-dim-penalty} \\
\midrule
\datasetname{Stream Benchmark 1} & 318464 (116480, 494848) & 368640 (163584, 546048) & 231936 (120576, 390912) \\
\datasetname{Stream Benchmark 2} & \textcolor{red}{154112 (88064, 183296)} & \textbf{82432 (56320, 92160)} & \textbf{82944 (54272, 95744)} \\
\datasetname{Random RBF Generator with Events} & \textcolor{red}{1099776 (421888, 1548288)} & \textbf{302080 (205312, 617728)} & \textbf{319488 (216832, 501504)} \\
\datasetname{Cassini} & 149504 (134656, 343296) & 147968 (109568, 175616) & 156672 (117504, 183296) \\
\datasetname{Chameleon DS3} & \textcolor{red}{734208 (604672, 2595584)} & \textcolor{red}{625152 (472832, 2160640)} & \textbf{278528 (274432, 636160)} \\
\datasetname{Chameleon DS4} & \textcolor{red}{695296 (587264, 2528768)} & \textcolor{red}{682496 (454400, 2352640)} & \textbf{251904 (246784, 288256)} \\
\datasetname{SCL-32} & \textbf{1164288 (34304, 1228288)} & - & \textcolor{red}{1272832 (1213952, 1360128)} \\
\datasetname{SCL-64} & \textbf{2300928 (60160, 2376448)} & - & \textcolor{red}{2385408 (2325504, 2529024)} \\
\bottomrule
\end{tabular}
}
\label{tab:dimpen_memory_max_bytes}
\end{table}

\begin{table}[htbp]
\centering
\caption{Statistical analysis for the Steady-State Memory Footprint [B].}
\resizebox{\textwidth}{!}{%
\begin{tabular}{l c c c }
\toprule
Dataset & \modelname{baseline} & \modelname{order-independent} & \modelname{no-dim-penalty} \\
\midrule
\datasetname{Stream Benchmark 1} & 67135 (17499, 322840) & 188480 (115516, 515336) & 187582 (110631, 497270) \\
\datasetname{Stream Benchmark 2} & 90504 (11791, 136409) & 120805 (35672, 210991) & 119442 (30617, 166601) \\
\datasetname{Random RBF Generator with Events} & 307390 (188323, 1218522) & 557951 (242417, 1366732) & 568623 (258952, 1366732) \\
\datasetname{Cassini} & 123210 (113289, 131625) & 102881 (49032, 118929) & 101585 (49032, 115658) \\
\datasetname{Chameleon DS3} & 317525 (314250, 489585) & 307577 (305326, 431567) & 309860 (303852, 431567) \\
\datasetname{Chameleon DS4} & 280885 (266411, 400677) & 264512 (256289, 395176) & 263681 (256289, 395176) \\
\datasetname{SCL-32} & \textbf{19476 (14216, 704242)} & \textbf{19596 (14637, 107708)} & \textcolor{red}{1200718 (1161079, 1366192)} \\
\datasetname{SCL-64} & \textbf{19579 (15685, 2292394)} & \textbf{18014 (15832, 100725)} & \textcolor{red}{2337305 (2308103, 2776593)} \\
\bottomrule
\end{tabular}
}
\label{tab:dimpen_memory_steadystate_bytes}
\end{table}

\begin{table}[htbp]
\centering
\caption{Statistical analysis for the Maximum Microcluster Count.}
\resizebox{\textwidth}{!}{%
\begin{tabular}{l c c c }
\toprule
Dataset & \modelname{baseline} & \modelname{order-independent} & \modelname{no-dim-penalty} \\
\midrule
\datasetname{Stream Benchmark 1} & 38 (15, 119) & 62 (41, 77) & 54 (33, 70) \\
\datasetname{Stream Benchmark 2} & 41 (16, 57) & 20 (14, 28) & 20 (15, 28) \\
\datasetname{Random RBF Generator with Events} & \textcolor{red}{228 (60, 430)} & \textbf{44 (28, 68)} & \textbf{45 (25, 70)} \\
\datasetname{Cassini} & 46 (42, 52) & 80 (34, 96) & 82 (34, 98) \\
\datasetname{Chameleon DS3} & 129 (123, 238) & 102 (93, 591) & 102 (93, 591) \\
\datasetname{Chameleon DS4} & 120 (110, 155) & 106 (99, 195) & 106 (99, 195) \\
\datasetname{SCL-32} & \textbf{39 (18, 84)} & - & \textcolor{red}{76 (54, 117)} \\
\datasetname{SCL-64} & \textbf{28 (24, 88)} & - & \textcolor{red}{83 (46, 215)} \\
\bottomrule
\end{tabular}
}
\label{tab:dimpen_mc_max}
\end{table}

\begin{table}[htbp]
\centering
\caption{Statistical analysis for the Steady-State Microcluster Count.}
\resizebox{\textwidth}{!}{%
\begin{tabular}{l c c c }
\toprule
Dataset & \modelname{baseline} & \modelname{order-independent} & \modelname{no-dim-penalty} \\
\midrule
\datasetname{Stream Benchmark 1} & 12 (3, 67) & 44 (28, 70) & 44 (21, 62) \\
\datasetname{Stream Benchmark 2} & 24 (6, 48) & 13 (8, 19) & 12 (8, 18) \\
\datasetname{Random RBF Generator with Events} & 68 (20, 289) & 34 (18, 46) & 31 (18, 54) \\
\datasetname{Cassini} & \textbf{36 (33, 44)} & \textcolor{red}{72 (24, 90)} & \textcolor{red}{73 (24, 94)} \\
\datasetname{Chameleon DS3} & 119 (114, 224) & 90 (78, 556) & 90 (79, 556) \\
\datasetname{Chameleon DS4} & 120 (110, 154) & 100 (95, 182) & 100 (95, 182) \\
\datasetname{SCL-32} & \textbf{8 (4, 39)} & - & \textcolor{red}{39 (16, 64)} \\
\datasetname{SCL-64} & \textbf{6 (5, 24)} & - & \textcolor{red}{56 (22, 211)} \\
\bottomrule
\end{tabular}
}
\label{tab:dimpen_mc_steadystate}
\end{table}


\section{Correction of the Typicality Formulation}
\label{sec:correction_typicality}

The evaluation of the original implementation reveals a critical mathematical discrepancy in the computation of a sample's membership to a macrocluster. The algorithm assigns a sample $x_k$ to a macrocluster by weighting the typicality of the sample against the active microclusters. However, the implemented typicality formulation evaluates an unnormalized intermediate expression, equivalent to $1 - \zeta^i(x_k)$. This expression contradicts the foundational definitions of the TEDA framework \cite{angelov2014outside}. In the formal mathematical model, absolute typicality is strictly defined as $\tau^i(x_k) = 1 - \xi^i(x_k)$. Furthermore, to construct valid density-based membership probabilities, the aggregation must utilize normalized typicality, defined as:

\begin{equation}
    t^i(x_k) = \frac{\tau^i(x_k)}{S^i_k - 2}
\end{equation}

By omitting the cardinality normalization factor $S^i_k - 2$, the original algorithm fails to appropriately scale the typicality across microclusters of varying sizes, inherently invalidating the resulting membership distributions. To restore theoretical consistency, the membership calculation is corrected to explicitly compute the exact normalized typicality $t^i(x_k)$. It is important to note that this normalized formulation is mathematically defined exclusively for microclusters containing more than two samples ($S^i_k > 2$). This domain constraint is naturally satisfied in the execution pipeline; because the macrocluster activation condition strictly filters out microclusters with $S^i_k \le 2$ prior to the membership evaluation, the denominator is mathematically guaranteed to be strictly positive.
